\chapter{Instrucciones de ejecución}
\label{ch:anexo}

\section{Compilación y ejecución}

El código se compila con Cargo, compilador de Rust. Desde la raíz del repositorio se ejecuta:
\begin{verbatim}
cargo build --release
\end{verbatim}

El binario principal se encuentra en \texttt{target/release/jsp} y recibe la ruta de la instancia,
un tiempo límite opcional y el subcomando de la estrategia.
Para el recocido simulado con los parámetros por defecto:
\begin{verbatim}
./target/release/jsp -f data/facil.txt simulated-annealing
./target/release/jsp -f data/medio.txt simulated-annealing
./target/release/jsp -f data/dificil.txt simulated-annealing
\end{verbatim}

Para ejecutar la experimentación completa de diez repeticiones por instancia:
\begin{verbatim}
./target/release/benchmark
\end{verbatim}

La ayuda de la interfaz de línea de comandos se obtiene con:
\begin{verbatim}
./target/release/jsp --help
./target/release/jsp simulated-annealing --help
\end{verbatim}

\section{Estructura del repositorio}

\begin{verbatim}
jsp/
|-- Cargo.toml
|-- README.txt
|-- data/
|   |-- facil.txt
|   |-- medio.txt
|   `-- dificil.txt
|-- src/
|   |-- domain/
|   |-- solver/
|   |-- ui/
|   `-- main.rs
|-- benches/
|   `-- benchmark.rs
`-- docs/
    `-- informe/
\end{verbatim}

El archivo \texttt{README.txt} contiene las instrucciones de ejecución y
la identificación de todos los integrantes del grupo,
conforme a lo solicitado en el enunciado de la tarea.

\section{Instancias utilizadas}

Las tres instancias se almacenan en la carpeta \texttt{data} y utilizan el formato estándar del \gls{jsp}:
la primera línea contiene el número de trabajos y de máquinas;
cada línea siguiente contiene pares (máquina, duración) en el orden tecnológico del trabajo correspondiente.

\begin{table}[htbp]
\centering
\caption{Instancias utilizadas en la experimentación.}
\label{tab:instancias}
\begin{tabular}{lrrr}
\toprule
Instancia & Trabajos & Máquinas & Operaciones \\
\midrule
Fácil   & 6  & 6  & 36 \\
Medio   & 15 & 10 & 150 \\
Difícil & 50 & 10 & 500 \\
\bottomrule
\end{tabular}
\end{table}
